\input{template.tex}

\ifPDFTeX
  \renewcommand{\rmdefault}{cmr}
\fi

\def\SheetNo{4}
\setlength{\emergencystretch}{2em}

\begin{document}
\sheettitle{Логика и комбинаторика.}
\begin{center}\textbf{Решения}\end{center}

\begin{problems}

\item Население острова состоит из $100$ рыцарей и $100$ лжецов.
Рыцари всегда правдивы, лжецы всегда лгут. Дружба взаимна, и у каждого
жителя есть хотя бы один друг. При опросе $100$ жителей утверждали,
что среди их друзей нет лжецов, а остальные $100$ - что среди их
друзей нет рыцарей. Каково наименьшее возможное число дружеских пар,
в которых один житель - рыцарь, а другой - лжец?
\par\smallskip\noindent\textbf{Решение.} Ответ: $50$. Пусть $a$ рыцарей сказали, что среди их друзей нет рыцарей. У каждого из них есть друг-лжец. Остальные $100-a$ человек, произнёсшие эту фразу, - лжецы; у каждого из них есть друг-рыцарь, иначе сказанное было бы правдой. Поэтому смешанных дружеских пар не меньше $\max(a,100-a)\ge50$.
\par Оценка достигается так. Разделим рыцарей и лжецов каждого типа на две группы по $50$ человек. Соединим попарно $50$ рыцарей с $50$ лжецами и не дадим этим жителям других друзей. Остальных рыцарей соединим между собой парами, и так же поступим с остальными лжецами. В смешанных парах рыцарь правдиво говорит, что у него нет друзей-рыцарей, а лжец лжёт теми же словами. Из остальных рыцари правдиво говорят, что у них нет друзей-лжецов, а лжецы лгут теми же словами. Получается ровно $100$ высказываний каждого вида и $50$ смешанных пар.\par

\item В волшебном городе есть неограниченное количество баров. Есть $5$ практиков, у которых есть $5$ предпочитаемых баров. У каждого из $m$ студентов тоже есть $5$ любимых баров в этом городе. При каком наибольшем $m$ каждый (студенты и практики) может пойти в один из своих любимых баров так, чтобы ни в одном баре не встретились практик и студент, вне зависимости от предпочтений? Требуется показать, что при большем $m$ можно подобрать предпочтения, при которых условие невыполнимо, а при найденном $m$ предложить способ для любых предпочтений.
\par\smallskip\noindent\textbf{Решение.} Ответ: $5^5-1=3124$. Если у двух практиков есть общий любимый бар, отправим обоих туда, а остальных троих - в любые любимые бары. Практики займут не более четырёх баров, поэтому у каждого студента останется хотя бы один незанятый любимый бар.
\par Если же пятёрки любимых баров практиков попарно не пересекаются, имеется ровно $5^5$ способов выбрать по одному бару для каждого практика. Пятёрка баров одного студента совпадает не более чем с одним таким выбором. При $m<5^5$ найдётся выбор, не совпадающий ни с одной студенческой пятёркой. Тогда каждый студент сможет пойти в свой любимый бар, свободный от практиков.
\par Для обратной оценки зададим практикам попарно непересекающиеся пятёрки баров, а $5^5$ студентам раздадим в качестве любимых все возможные выборы по одному бару из каждой пятёрки. Какие бы бары ни выбрали практики, одному студенту достанется в точности множество занятых ими баров. Он не сможет пойти ни в один из них.\par

\item Являются ли полными следующие системы булевых функций?
\begin{parts}
  \item $\{\neg,\vee\}$;
  \item $\{\oplus,\vee\}$;
  \item $\{\oplus,\supset\}$.
\end{parts}
\par\smallskip\noindent\textbf{Решение.} \textbf{а)} Да: $x\mathbin{\&}y=\neg(\neg x\vee\neg y)$, а отрицания, конъюнкции и дизъюнкции достаточно для записи любой СДНФ.
\par\textbf{б)} Нет: обе операции сохраняют нуль, поэтому любая выраженная через них функция равна нулю на наборе из одних нулей. Например, константу $1$ выразить нельзя.
\par\textbf{в)} Да: $1=x\supset x$, затем $\neg x=x\oplus1$ и $x\mathbin{\&}y=\neg(x\supset\neg y)$. Получен полный набор операций.\par

\item Докажите тождество
\[
  C_n^0+C_{n+1}^1+C_{n+2}^2+\dots+C_{n+k}^k=C_{n+k+1}^k.
\]
\par\smallskip\noindent\textbf{Решение.} По формуле Паскаля
\[
C_{n+j}^{j}=C_{n+j+1}^{j}-C_{n+j}^{j-1}.
\]
Просуммируем эти равенства по $j=0,\ldots,k$, полагая $C_n^{-1}=0$. Промежуточные слагаемые сокращаются, остаётся $C_{n+k+1}^{k}$.\par

\item Сто мудрецов ставят в ряд и каждому надевают либо красный,
либо синий колпак; колпаков каждого цвета имеется по $100$.
Мудрец видит только колпаки стоящих перед ним. Затем мудрецы по одному,
в выбранном ими порядке, объявляют предполагаемый цвет собственного
колпака. Все слышат предыдущие ответы; других сигналов подавать нельзя.
Тот, кто ошибся, будет казнён. До начала испытания мудрецы могут
согласовать стратегию. Какое максимальное число спасённых они могут
обеспечить при любом распределении колпаков?
\par\smallskip\noindent\textbf{Решение.} Ответ: $99$. Пусть мудрецы говорят от последнего в ряду к первому. Первый называет красный цвет, если число красных колпаков впереди нечётно, и синий, если оно чётно. Его ответ передаёт остальным чётность числа красных колпаков на них. Каждый следующий мудрец знает цвета колпаков впереди и из уже верных ответов - цвета колпаков позади, поэтому однозначно восстанавливает свой цвет.
\par Спасти гарантированно всех нельзя: первый отвечающий ещё не слышал ответов и не видит собственного колпака. При неизменных видимых ему колпаках его собственный цвет можно поменять, и в одном из двух случаев его ответ будет ошибочным.\par

\item Между некоторыми жителями острова установлены взаимные дружеские
отношения. Каждый житель - либо неизменно правдивый рыцарь, либо
неизменно лгущий лжец. Все жители утверждают, что рыцари составляют
больше половины их друзей. Докажите, что среди всех дружеских пар
не меньше половины составляют пары из двух рыцарей или двух лжецов.
\par\smallskip\noindent\textbf{Решение.} Обозначим через $e_{RR},e_{LL},e_{RL}$ количества дружеских пар соответствующих типов. У каждого рыцаря друзей-рыцарей больше, чем друзей-лжецов. Суммируя по рыцарям, получаем $2e_{RR}\ge e_{RL}$. У каждого лжеца друзей-рыцарей не больше половины всех друзей, поскольку его утверждение ложно. Значит, друзей-лжецов у него не меньше, чем друзей-рыцарей; суммирование даёт $2e_{LL}\ge e_{RL}$. Следовательно, $e_{RR}+e_{LL}\ge e_{RL}$, то есть однотипных пар не меньше половины.\par

\item Курс лекций состоял из $35$ пар. На каждой паре присутствовало ровно $27$ студентов. Оказалось, что на любых трёх занятиях был ровно один общий студент. Верно ли, что был студент, посетивший все пары?
\par\smallskip\noindent\textbf{Решение.} Да. Предположим, что никто не посетил все пары, и выберем студента, посетившего максимальное число $m$ пар. Тогда $3\le m<35$. Возьмём пару $A$, которую он пропустил. Каждой неупорядоченной паре $X,Y$ из посещённых им занятий соответствует единственный студент, присутствовавший на $A,X,Y$. Для разных пар $\{X,Y\}$ эти студенты различны: иначе один из них посетил бы хотя бы три занятия из посещённых выбранным студентом, и у этой тройки было бы два общих студента. Поэтому на занятии $A$ присутствовало не меньше $C_m^2$ студентов. Получаем $C_m^2\le27$, откуда $m\le7$.
\par Пусть $S$ - множество студентов, посетивших хотя бы три пары, $s=|S|$, а $a_i$ - число посещений $i$-го из них. Тогда $3\le a_i\le7$, а каждая тройка занятий учитывается ровно у одного студента. Следовательно,
\[
\sum_{i=1}^{s}C_{a_i}^{3}=C_{35}^{3}=6545.
\]
Поскольку $C_{a_i}^{3}\le C_7^3=35$, имеем $s\ge6545/35=187$. С другой стороны,
\[
C_{a_i}^{3}=\frac{a_i(a_i-1)(a_i-2)}{6}\le7(a_i-2).
\]
Всего посещений $35\cdot27=945$, поэтому $\sum_i a_i\le945$ и
\[
6545\le7\left(\sum_{i=1}^{s}a_i-2s\right)\le7(945-2s).
\]
Отсюда $935\le945-2s$, то есть $s\le5$, что противоречит $s\ge187$. Значит, некоторый студент посетил все $35$ пар.\par

\item Докажите формулу (а) алгебраически и (б) комбинаторно:
\[
  C_n^k C_k^l = C_n^l C_{n-l}^{k-l}.
\]
\par\smallskip\noindent\textbf{Решение.} \textbf{а)} При $0\le l\le k\le n$ обе части равны
\[
\frac{n!}{l!\,(k-l)!\,(n-k)!}.
\]
\textbf{б)} Обе части считают пары подмножеств $(L,K)$, для которых $L\subseteq K\subseteq\{1,\ldots,n\}$, $|L|=l$ и $|K|=k$. Слева сначала выбирается $K$, затем $L\subseteq K$; справа сначала $L$, затем остальные $k-l$ элементов множества $K$.\par

\item Среди трёх подозреваемых - Брауна, Джонса и Смита - ровно один
совершил преступление. Каждый сообщил следователю два утверждения.
Браун отрицает как свою причастность, так и причастность Джонса.
Смит отрицает свою причастность и обвиняет Брауна.
Джонс отрицает причастность Брауна и обвиняет Смита.
Известно, что у одного подозреваемого оба утверждения истинны,
у другого оба ложны, а у третьего одно истинно и одно ложно.
Определите виновного.
\par\smallskip\noindent\textbf{Решение.} Виновен Браун. Если виновен он, высказывания Брауна имеют значения «ложь, правда», Смита - «правда, правда», Джонса - «ложь, ложь», как и требуется. Если виновен Джонс, у каждого подозреваемого одно истинное и одно ложное утверждение. Если виновен Смит, у Брауна и Джонса оба утверждения истинны, а у Смита оба ложны. Последние два случая не подходят.\par

\item На острове рыцари никогда не лгут, а лжецы никогда не говорят правды.
В хороводе из $21$ жителя представлены оба типа островитян.
Каждый участник утверждает, что его два соседа принадлежат к разным типам:
один из них рыцарь, другой - лжец. Определите число рыцарей в хороводе.
\par\smallskip\noindent\textbf{Решение.} Ответ: $14$. Обозначим через $x_i\in\{0,1\}$ тип $i$-го жителя: $1$ для рыцаря и $0$ для лжеца. Высказывание истинно тогда и только тогда, когда типы соседей различны, поэтому по модулю $2$
\[
x_i=x_{i-1}+x_{i+1},\qquad x_{i+1}=x_i+x_{i-1}.
\]
Отсюда $x_{i+2}=x_{i-1}$: последовательность типов имеет период $3$. В каждой тройке выполняется $x_{i-1}+x_i+x_{i+1}=0\pmod2$, поэтому в ней либо все нули, либо ровно две единицы. Оба типа встречаются, значит, каждая из семи троек содержит двух рыцарей.\par

\item Не пользуясь таблицами истинности, приведите к СКНФ и СДНФ следующие функции:
\begin{parts}
  \item $((P\supset Q)\supset(R\supset\neg P))\supset(\neg Q\supset\neg R)$;
  \item $(((((P\supset Q)\supset\neg P)\supset\neg Q)\supset\neg R)\supset R)$.
\end{parts}
\par\smallskip\noindent\textbf{Решение.} \textbf{а)} Последовательно раскрывая импликации, получаем
\[
(P\supset Q)\supset(R\supset\neg P)
\equiv\neg P\vee\neg Q\vee\neg R,
\qquad
\neg Q\supset\neg R\equiv Q\vee\neg R.
\]
Поэтому вся функция равна $Q\vee\neg R$. Её СКНФ:
\[
(P\vee Q\vee\neg R)\mathbin{\&}(\neg P\vee Q\vee\neg R).
\]
Её СДНФ:
\[
\begin{aligned}
&(\neg P\mathbin{\&}\neg Q\mathbin{\&}\neg R)
\vee(\neg P\mathbin{\&}Q\mathbin{\&}\neg R)
\vee(\neg P\mathbin{\&}Q\mathbin{\&}R)\\
{}\vee{}&(P\mathbin{\&}\neg Q\mathbin{\&}\neg R)
\vee(P\mathbin{\&}Q\mathbin{\&}\neg R)
\vee(P\mathbin{\&}Q\mathbin{\&}R).
\end{aligned}
\]
\textbf{б)} Упрощение внутренних частей даёт
\[
((P\supset Q)\supset\neg P)\equiv\neg P\vee\neg Q,
\qquad
((\neg P\vee\neg Q)\supset\neg Q)\equiv P\vee\neg Q.
\]
Следовательно, исходная функция равна $((P\vee\neg Q)\supset\neg R)\supset R\equiv R$. Её СКНФ:
\[
\begin{aligned}
&(P\vee Q\vee R)\mathbin{\&}(P\vee\neg Q\vee R)\\
{}\mathbin{\&}{}&(\neg P\vee Q\vee R)\mathbin{\&}(\neg P\vee\neg Q\vee R).
\end{aligned}
\]
Её СДНФ:
\[
\begin{aligned}
&(\neg P\mathbin{\&}\neg Q\mathbin{\&}R)
\vee(\neg P\mathbin{\&}Q\mathbin{\&}R)\\
{}\vee{}&(P\mathbin{\&}\neg Q\mathbin{\&}R)
\vee(P\mathbin{\&}Q\mathbin{\&}R).
\end{aligned}
\]\par

\item На острове $2023$ жителя, каждый из которых - либо рыцарь,
всегда отвечающий правдиво, либо лжец, всегда дающий ложный ответ.
Антрополог спросил каждого жителя о числе его друзей на острове;
дружба взаимна. Полученные ответы - целые числа от $0$ до $2022$,
причём одинаковые ответы допускаются. По результатам опроса антрополог
смог доказать, что число лжецов не меньше $k$. Найдите наибольшее $k$,
для которого такая ситуация возможна.
\par\smallskip\noindent\textbf{Решение.} Ответ: $1011$. Разделим возможные ответы на малые ($0,\ldots,1011$) и большие ($1012,\ldots,2022$). Как минимум $1012$ жителей дали ответы одной из этих групп. Если это малые ответы, выберем $1012$ таких жителей, а остальных $1011$ объявим вспомогательными. Каждому выбранному жителю с ответом $d\le1011$ проведём ровно $d$ дружеских связей со вспомогательными жителями; других связей не нужно. Тогда все $1012$ выбранных ответили правдиво.
\par Если выбранные ответы большие, для каждого из них положим $d'=2022-d\le1010$, построим вспомогательный граф предыдущим способом для степеней $d'$ и заменим его дополнением. В итоговом графе у выбранного жителя будет ровно $d$ друзей. Значит, при любых ответах возможно не менее $1012$ рыцарей, и доказать наличие более $1011$ лжецов нельзя.
\par Для достижения оценки пусть $1011$ жителей назовут $0$, а $1012$ - $2022$. Житель с нулём друзей и житель, друживший со всеми, не могут оба говорить правду. Поэтому все правдивые ответы принадлежат лишь одной из групп; их не больше $1012$. Следовательно, лжецов не меньше $1011$.\par

\item
\begin{enumerate}
  \item[а)] В ряд выстроены $n$ школьников. Сколькими способами можно выбрать из них $k$ дежурных так, чтобы никакие два дежурных не стояли рядом?
  \item[б)] За круглым столом короля Артура сидят $n$ рыцарей. Каждый из них враждует со своими соседями. Король хочет составить отряд из $k$ рыцарей. Сколькими способами это можно сделать так, чтобы в этом отряде не было врагов?
\end{enumerate}
\par\smallskip\noindent\textbf{Решение.} \textbf{а)} Пусть выбранные места имеют номера $i_1<\dots<i_k$. Тогда числа $j_s=i_s-(s-1)$ образуют произвольный возрастающий набор из $\{1,\ldots,n-k+1\}$. Ответ: $C_{n-k+1}^k$ при $0\le k\le\lceil n/2\rceil$ и $0$ при больших $k$.
\par\textbf{б)} При $k=0$ ответ равен $1$, а при $k>n/2$ - $0$. Если $1\le k\le n/2$, зафиксируем одного выбранного рыцаря и рассмотрим длины $k$ промежутков до следующих выбранных по кругу. Все длины положительны, а их сумма равна $n-k$; наборов длин $C_{n-k-1}^{k-1}$. Выбрать начального рыцаря можно $n$ способами, но каждый отряд посчитан $k$ раз. Ответ:
\[
\frac{n}{k}C_{n-k-1}^{k-1}.
\]\par

\item Докажите, что число $\displaystyle\frac{(n^2)!}{(n!)^{n+1}}$ является целым для натуральных значений $n$.
\par\smallskip\noindent\textbf{Решение.} Рассмотрим разбиения множества из $n^2$ различных элементов на $n$ неупорядоченных групп по $n$ элементов. Если группы временно пронумеровать, разбиений получится $(n^2)!/(n!)^n$. Каждое неупорядоченное разбиение при этом посчитано $n!$ раз. Следовательно, искомое число равно количеству таких разбиений и потому целое.\par

\item Найдите количество способов, которыми два человека могут разделить $2n$ предметов одного сорта, $2n$ предметов второго сорта и $2n$ предметов третьего сорта так, чтобы каждый получил $3n$ предметов.
\par\smallskip\noindent\textbf{Решение.} Пусть первый человек получает $x,y,z$ предметов соответствующих сортов. Тогда распределение однозначно задаётся решением
\[
x+y+z=3n,\qquad 0\le x,y,z\le2n
\]
в целых числах; второй человек получает оставшиеся предметы. Без верхних ограничений число неотрицательных решений равно $C_{3n+2}^{2}$. Вычтем решения, в которых хотя бы одна переменная превышает $2n$.
\par Для случая $x\ge2n+1$ положим $w=x-(2n+1)$. Получим уравнение $y+z+w=n-1$ в неотрицательных целых числах, у которого $C_{n+1}^{2}$ решений. Для $y$ и $z$ подсчёт аналогичен. Эти три случая не пересекаются: две переменные, превышающие $2n$, дали бы сумму не меньше $4n+2>3n$. Поэтому ответ равен
\[
C_{3n+2}^{2}-3C_{n+1}^{2}=3n^2+3n+1.
\]\par

\item В Министерстве Пропаганды работали честные сотрудники (которые всегда говорили правду) и патологические лжецы (которые всегда лгали). Каждый сотрудник знал про любого другого, честный тот или лжец. Вновь назначенный Министр Пропаганды у каждого подчинённого спросил про какого-то другого подчинённого, честный тот или лжец. При этом ни про какого подчинённого Министр не спрашивал дважды. После опроса Министр выгнал из Министерства всех сотрудников, которые были названы честными. Могло ли у Министра остаться ровно $2013$ подчинённых?
\par\smallskip\noindent\textbf{Решение.} Нет. Каждый сотрудник ответил на один вопрос, и о каждом спросили не более одного раза. Значит, направленные связи «отвечающий $\to$ тот, о ком спросили» разбивают сотрудников на непересекающиеся циклы. Сотрудник называет другого честным тогда и только тогда, когда они одного типа: честный говорит правду, а лжец меняет ответ на противоположный. Поэтому остаются ровно те сотрудники, тип которых отличается от типа следующего человека в их цикле. В каждом замкнутом цикле число таких смен типа чётно. Сумма по всем циклам также чётна, а $2013$ нечётно.\par

\item
\begin{enumerate}
  \item[а)] (Формула Вандермонда.) Пусть $m,n,k\in\mathbb N$, причём $m,n\geq k$. Докажите, что
  \[
    C_n^0C_m^k+C_n^1C_m^{k-1}+\dots+C_n^kC_m^0=C_{m+n}^k.
  \]
  \item[б)] Докажите алгебраически, что
  \[
    (C_n^0)^2+(C_n^1)^2+\dots+(C_n^n)^2=C_{2n}^n.
  \]
\end{enumerate}
\par\smallskip\noindent\textbf{Решение.} \textbf{а)} В разложении $(1+x)^n(1+x)^m$ коэффициент при $x^k$ равен $\sum_{j=0}^k C_n^jC_m^{k-j}$. Но произведение равно $(1+x)^{n+m}$, где тот же коэффициент есть $C_{n+m}^k$.
\par\textbf{б)} Подставим $m=n$, $k=n$ в пункт а) и воспользуемся симметрией $C_n^{n-j}=C_n^j$.\par

\item Докажите тождества (во втором пункте $n\geq2$):
\begin{enumerate}
  \item[а)] $C_n^1+2C_n^2+\dots+nC_n^n=n2^{n-1}$;
  \item[б)] $C_n^1-2C_n^2+\dots+(-1)^{n-1}nC_n^n=0$.
\end{enumerate}
\par\smallskip\noindent\textbf{Решение.} Продифференцируем биномиальное тождество $(1+x)^n=\sum_{k=0}^n C_n^k x^k$:
\[
n(1+x)^{n-1}=\sum_{k=1}^n kC_n^k x^{k-1}.
\]
При $x=1$ получаем пункт а), а при $x=-1$ и $n\ge2$ - пункт б).\par

\item Докажите, что булеву функцию $f(x_1,\ldots,x_n)$ можно выразить через $\mathbin{\&}$, $\vee$, $\supset$ тогда и только тогда, когда в её СКНФ отсутствует дизъюнкт $(\neg x_1\vee\cdots\vee\neg x_n)$.
\par\smallskip\noindent\textbf{Решение.} Все три разрешённые операции принимают значение $1$ на единичных аргументах. Поэтому любая построенная из них функция удовлетворяет $f(1,\ldots,1)=1$. В СКНФ дизъюнкт $(\neg x_1\vee\cdots\vee\neg x_n)$ присутствует в точности тогда, когда $f(1,\ldots,1)=0$.
\par Пусть теперь этого дизъюнкта нет. Каждый дизъюнкт СКНФ содержит хотя бы одну переменную без отрицания и имеет вид
\[
\neg x_{i_1}\vee\cdots\vee\neg x_{i_r}\vee x_{j_1}\vee\cdots\vee x_{j_s},\qquad s\ge1.
\]
Он равносилен $(x_{i_1}\mathbin{\&}\cdots\mathbin{\&}x_{i_r})\supset(x_{j_1}\vee\cdots\vee x_{j_s})$. При $r=0$ остаётся дизъюнкция положительных переменных. Соединив полученные выражения конъюнкцией, получаем $f$. Если СКНФ пуста, то $f\equiv1$ выражается формулой $x_1\supset x_1$.\par

\item Существует ли полная система, состоящая из одной булевой функции?
\par\smallskip\noindent\textbf{Решение.} Да. Возьмём штрих Шеффера $x\mid y=\neg(x\mathbin{\&}y)$. Из него выражаются отрицание $\neg x=x\mid x$ и конъюнкция $x\mathbin{\&}y=(x\mid y)\mid(x\mid y)$. Этих двух операций достаточно для построения любой булевой функции по её СДНФ.\par

\item Сколько имеется четырёхзначных чисел, у которых каждая следующая цифра меньше предыдущей?
\par\smallskip\noindent\textbf{Решение.} Ответ: $C_{10}^{4}=210$. Нужно выбрать четыре различные цифры из $0,\ldots,9$; их убывающий порядок задаётся однозначно. Первая цифра в таком порядке - наибольшая из четырёх, поэтому она не равна нулю.\par

\item Дано число $53\underbrace{00\ldots00}_{100\text{ нулей}}35$. Требуется заменить некоторые два нуля на ненулевые цифры так, чтобы получилось число, делящееся на $495$. Сколькими способами это можно сделать?
\par\smallskip\noindent\textbf{Решение.} Ответ: $22100$. Поскольку $495=5\cdot9\cdot11$, проверим делимость на каждый из этих взаимно простых множителей. Последняя цифра остаётся равной $5$, поэтому делимость на $5$ выполнена. Сумма исходных цифр равна $16$, а знакочередующаяся сумма - нулю. Если вставлены цифры $a,b\in\{1,\ldots,9\}$, то для делимости на $9$ необходимо и достаточно $16+a+b\equiv0\pmod9$.
\par Среди ста нулей по $50$ позиций каждой чётности. Если выбранные позиции одной чётности, условие делимости на $11$ имеет вид $a+b\equiv0\pmod{11}$. Поэтому $a+b=11$; это также обеспечивает делимость на $9$. Для каждой пары позиций есть восемь вариантов цифр: $(2,9),(3,8),\ldots,(9,2)$, где первой считается цифра на левой позиции. Пар позиций $2C_{50}^{2}=50\cdot49$, так что этот случай даёт $50\cdot49\cdot8=19600$ способов.
\par Если позиции разной чётности, то $a-b\equiv0\pmod{11}$, откуда $a=b$. Условие делимости на $9$ принимает вид $16+2a\equiv0\pmod9$; среди ненулевых цифр ему удовлетворяет только $a=1$. Пар позиций $50\cdot50=2500$, и для каждой подходит единственная пара цифр $(1,1)$. Итого получаем $19600+2500=22100$ способов.\par

\item Вычислите сумму $\displaystyle\sum_{k=0}^{n} 2^k C_n^k$.
\par\smallskip\noindent\textbf{Решение.} По формуле бинома Ньютона
\[
\sum_{k=0}^{n}2^k C_n^k=(1+2)^n=3^n.
\]\par

\item Докажите, что набор простых дизъюнктов выполним (то есть можно так подобрать значения переменных, чтобы все эти дизъюнкты стали истинными) тогда и только тогда, когда из них при помощи правила резолюции нельзя получить пустой дизъюнкт.
\par\smallskip\noindent\textbf{Решение.} Правило резолюции из дизъюнктов $A\vee x$ и $B\vee\neg x$ выводит $A\vee B$. При любой оценке, удовлетворяющей обоим исходным дизъюнктам, истинен и вывод: если бы $A$ и $B$ были ложны, пришлось бы одновременно положить $x=1$ и $x=0$. Поэтому из выполнимого набора пустой дизъюнкт вывести нельзя.
\par Обратно, устраним одну переменную $x$. Оставим все дизъюнкты без $x$ и добавим все резольвенты пар, содержащих соответственно $x$ и $\neg x$. Полученный набор выполним тогда и только тогда, когда исходный. В одну сторону это следует из корректности резолюции. В другую: если оценка оставшихся переменных удовлетворяет всем резольвентам, то она не может одновременно требовать $x=1$ в одном исходном дизъюнкте и $x=0$ в другом; значит, её можно продолжить до удовлетворяющей оценки переменной $x$. Повторяя устранение для всех переменных, в случае невыполнимости придём к пустому дизъюнкту. Каждый добавленный на пути дизъюнкт получен резолюцией, что и требовалось.\par

\item Каждый житель острова людоедов принадлежит к одному из двух племён - рыцарей, которые всегда говорят правду, или лжецов, которые всегда лгут. Однажды $2013$ островитян встали в круг и каждый заявил: «Оба моих соседа - лжецы». Тут началась перебранка, и одного из островитян съели. После этого оставшиеся $2012$ островитян снова встали в круг, возможно, в другом порядке, и каждый заявил: «Оба моих соседа не из моего племени». Кого съели - рыцаря или лжеца?
\par\smallskip\noindent\textbf{Решение.} Съели рыцаря. В первом круге каждый рыцарь окружён лжецами. Между соседними рыцарями может стоять не более двух лжецов: при трёх лжецах подряд средний говорил бы правду. Следовательно, лжецов не больше удвоенного числа рыцарей, и среди $2013$ жителей было не меньше $\lceil2013/3\rceil=671$ рыцаря.
\par Во втором круге каждый рыцарь вновь окружён лжецами. Между соседними рыцарями не может стоять ровно один лжец: его соседи тогда были бы не из его племени, и он сказал бы правду. Значит, после удаления между рыцарями стоит как минимум по два лжеца. Среди $2012$ жителей рыцарей не больше $\lfloor2012/3\rfloor=670$. Их число уменьшилось, поэтому удалённый житель был рыцарем.\par

\end{problems}
\end{document}
